\documentclass[10pt,letterpaper]{article}
\usepackage[utf8]{inputenc}
\usepackage[intlimits]{amsmath}
\usepackage{amsfonts}
\usepackage{amssymb}
\usepackage{ragged2e}
\usepackage[letterpaper, margin=0.5in]{geometry}
\usepackage{graphicx}
\usepackage{cancel}
\usepackage{mathtools}
\usepackage{tabularx}
\usepackage{arydshln}
\usepackage{tensor}
\usepackage{array}
\usepackage{xcolor}
\usepackage[boxed]{algorithm}
\usepackage[noend]{algpseudocode}
\usepackage{listings}
\usepackage{textcomp}
% \usepackage[pdf,tmpdir,singlefile]{graphviz}
\usepackage{mathrsfs}
\usepackage{bbm}
\usepackage{tikz}
\usepackage{tikz-cd}
\usepackage{enumitem}
\usepackage{arydshln}
\usepackage{relsize}
\usepackage{multicol}
\usepackage{scalerel}
\usepackage{upgreek}
\usepackage{ifthen}

\usetikzlibrary{bayesnet}
\setlist{noitemsep}

%%%%%%%%%%%%%%%%%%%%%%%%%%%%%
% Formatting commands
%%%%%%%%%%%%%%%%%%%%%%%%%%%%%
\newcommand{\n}{\hfill\break}
\newcommand{\nn}{\vspace{0.5\baselineskip}\n}
\newcommand{\up}{\vspace{-\baselineskip}}
\newcommand{\hangblock}[2]{\par\noindent\settowidth{\hangindent}{\textbf{#1: }}\textbf{#1: }\nolinebreak#2}
\newcommand{\lemma}[1]{\hangblock{Lemma}{#1}}
\newcommand{\defn}[1]{\hangblock{Defn}{#1}}
\newcommand{\thm}[1]{\hangblock{Thm}{#1}}
\newcommand{\cor}[1]{\hangblock{Cor}{#1}}
\newcommand{\prop}[1]{\hangblock{Prop}{#1}}
\newcommand{\ex}[1]{\hangblock{Ex}{#1}}
\newcommand{\exer}[1]{\hangblock{Exer}{#1}}
\newcommand{\fact}[1]{\hangblock{Fact}{#1}}
\newcommand{\remark}[1]{\hangblock{Remark}{#1}}
\newcommand{\proven}{\;$\square$\n}
\newcommand{\problem}[1]{\par\noindent{\nolinebreak#1}\n}
\newcommand{\problempart}[2]{\par\noindent\indent{}\settowidth{\hangindent}{\textbf{(#1)} \indent{}}\textbf{(#1) }\nolinebreak#2\n}
\newcommand{\ptxt}[1]{\textrm{\textnormal{#1}}}
\newcommand{\inlineeq}[1]{\centerline{$\displaystyle #1$}}
\newcommand{\pageline}{\up\par\noindent\rule{\textwidth}{0.1pt}}

%%%%%%%%%%%%%%%%%%%%%%%%%%%%%
% Math commands
%%%%%%%%%%%%%%%%%%%%%%%%%%%%%
% Set Theory
\newcommand{\card}[1]{\left|#1\right|}
\newcommand{\set}[1]{\left\{#1\right\}}
\newcommand{\setmid}{\;\middle|\;}
\newcommand{\ps}[1]{\mathcal{P}\left(#1\right)}
\newcommand{\pfinite}[1]{\mathcal{P}^{\ptxt{finite}}\left(#1\right)}
\newcommand{\naturals}{\mathbb{N}}
\newcommand{\N}{\naturals}
\newcommand{\integers}{\mathbb{Z}}
\newcommand{\Z}{\integers}
\newcommand{\rationals}{\mathbb{Q}}
\newcommand{\Q}{\rationals}
\newcommand{\reals}{\mathbb{R}}
\newcommand{\R}{\reals}
\newcommand{\complex}{\mathbb{C}}
\newcommand{\C}{\complex}
\newcommand{\halfPlane}{\mathbb{H}}
\let\HSym\H
\let\H\relax
\newcommand{\H}{\halfPlane}
\newcommand{\comp}{^{\complement}}
\DeclareMathOperator{\Hom}{Hom}
\newcommand{\Ind}{\mathbbm{1}}
\newcommand{\cut}{\setminus}
\DeclareMathOperator{\elem}{elem}

% Differentiable Manifolds
\newcommand{\RP}{\mathbb{RP}}
\newcommand{\CP}{\mathbb{RP}}
\newcommand{\osubset}{\overset{\mathclap{\scalebox{0.5}{\ptxt{open}}}}{\subset}}
\newcommand{\osubseteq}{\overset{\mathclap{\scalebox{0.5}{\ptxt{open}}}}{\subseteq}}
\newcommand{\osupset}{\overset{\mathclap{\scalebox{0.5}{\ptxt{open}}}}{\supset}}
\newcommand{\osupseteq}{\overset{\mathclap{\scalebox{0.5}{\ptxt{open}}}}{\supseteq}}
\newcommand{\pdat}[3]{\left.\pd{#1}{#2}\right|_{#3}}
\DeclareMathOperator{\so}{so}
\DeclareMathOperator{\codim}{codim}
\DeclareMathOperator{\Diff}{Diff}
\let\dSym\d
\let\d\relax
\newcommand{\d}{\partial}
\DeclareMathOperator{\gl}{gl}

% Graph Theory
\let\deg\relax
\DeclareMathOperator{\deg}{deg}
\newcommand{\degp}{\ptxt{deg}^{+}}
\newcommand{\degn}{\ptxt{deg}^{-}}
\newcommand{\precdot}{\mathrel{\ooalign{$\prec$\cr\hidewidth\hbox{$\cdot\mkern0.5mu$}\cr}}}
\newcommand{\succdot}{\mathrel{\ooalign{$\cdot\mkern0.5mu$\cr\hidewidth\hbox{$\succ$}\cr\phantom{$\succ$}}}}
\DeclareMathOperator{\cl}{cl}
\DeclareMathOperator{\affdim}{affdim}

% Probability
\newcommand{\parSymbol}{\P}
\newcommand{\Prob}{\mathbb{P}}
\renewcommand{\P}{\Prob}
\newcommand{\Avg}{\mathbb{E}}
\newcommand{\E}{\Avg}
\DeclareMathOperator{\Var}{Var}
\DeclareMathOperator{\cov}{cov}
\DeclareMathOperator{\Unif}{Unif}
\DeclareMathOperator{\Binom}{Binom}
\newcommand{\CI}{\mathrel{\text{\scalebox{1.07}{$\perp\mkern-10mu\perp$}}}}

% Standard Math
\newcommand{\inv}{^{-1}}
\newcommand{\abs}[1]{\left|#1\right|}
\newcommand{\ceil}[1]{\left\lceil{}#1\right\rceil{}}
\newcommand{\floor}[1]{\left\lfloor{}#1\right\rfloor{}}
\newcommand{\conj}[1]{\overline{#1}}
\newcommand{\of}{\circ}
\newcommand{\tri}{\triangle}
\newcommand{\inj}{\hookrightarrow}
\newcommand{\surj}{\twoheadrightarrow}
\newcommand{\ndiv}{\nmid}
\renewcommand{\epsilon}{\varepsilon}
\newcommand{\divides}{\mid}
\newcommand{\ndivides}{\nmid}
\DeclareMathOperator{\lcm}{lcm}
\DeclareMathOperator{\sgn}{sgn}
\newcommand{\map}[4]{\!\!\!\begin{array}[t]{rcl}#1 & \!\!\!\!\to & \!\!\!\!#2\\ {}#3 & \!\!\!\!\mapsto & \!\!\!\!#4\end{array}}
\newcommand{\bigsum}[2]{\smashoperator[lr]{\sum_{\scalebox{#1}{$#2$}}}}
\DeclareMathOperator{\gcf}{gcf}
\newcommand{\restr}[2]{\left.#1\right|_{#2}}

% Linear Algebra
\newcommand{\Id}{\textrm{\textnormal{Id}}}
\newcommand{\im}{\textrm{\textnormal{im}}}
\newcommand{\norm}[1]{\abs{\abs{#1}}}
\newcommand{\tpose}{^{T}\!}
\newcommand{\iprod}[1]{\left<#1\right>}
\newcommand{\giprod}{\iprod{\;\,,\;}}
\DeclareMathOperator{\tr}{tr}
\DeclareMathOperator{\trace}{tr}
\newcommand{\chgBasMat}[3]{\!\!\tensor*[_{#1}]{\left[#2\right]}{_{#3}}}
\newcommand{\vecBas}[2]{\tensor*[]{\left[#1\right]}{_{#2}}}
\DeclareMathOperator{\GL}{GL}
\DeclareMathOperator{\Mat}{Mat}
\DeclareMathOperator{\vspan}{span}
\DeclareMathOperator{\rank}{rank}
\newcommand{\V}[1]{\vec{#1}}
\DeclareMathOperator{\proj}{proj}
\DeclareMathOperator{\compProj}{comp}
\DeclareMathOperator{\row}{row}
\newcommand{\smallPMatrix}[1]{\paren{\begin{smallmatrix}#1\end{smallmatrix}}}
\newcommand{\smallBMatrix}[1]{\brack{\begin{smallmatrix}#1\end{smallmatrix}}}
\newcommand{\pmat}[1]{\begin{pmatrix}#1\end{pmatrix}}
\newcommand{\bmat}[1]{\begin{bmatrix}#1\end{bmatrix}}
\newcommand{\dual}{^{*}}
\newcommand{\pinv}{^{\dagger}}
\newcommand{\horizontalMatrixLine}{\ptxt{\rotatebox[origin=c]{-90}{$|$}}}
\DeclareMathOperator{\range}{range}
\DeclareMathOperator{\Symm}{Symm}
\DeclareMathOperator{\SU}{SU}

% Multilinear Algebra
\let\Lsym\L
\let\L\relax
\DeclareMathOperator{\L}{\mathscr{L}}
\DeclareMathOperator{\A}{\mathcal{A}}
\DeclareMathOperator{\Alt}{Alt}
\DeclareMathOperator{\Sym}{Sym}
\newcommand{\ot}{\otimes}
\newcommand{\ox}{\otimes}
\DeclareMathOperator{\asc}{asc}
\DeclareMathOperator{\asSet}{set}
\DeclareMathOperator{\sort}{sort}
\DeclareMathOperator{\ringA}{\mathring{A}}
\DeclareMathOperator{\Sh}{Sh}

% Topology
\newcommand{\closure}[1]{\overline{#1}}
\newcommand{\uball}{\mathcal{U}}
\DeclareMathOperator{\Int}{Int}
\DeclareMathOperator{\Ext}{Ext}
\DeclareMathOperator{\Bd}{Bd}
\DeclareMathOperator{\rInt}{rInt}
\DeclareMathOperator{\ch}{ch}
\DeclareMathOperator{\ah}{ah}
\newcommand{\LargerTau}{\mathlarger{\mathlarger{\mathlarger{\mathlarger{\tau}}}}}
\newcommand{\Tau}{\mathcal{T}}

% Analysis
\DeclareMathOperator{\Graph}{Graph}
\DeclareMathOperator{\epi}{epi}
\DeclareMathOperator{\hypo}{hypo}
\DeclareMathOperator{\supp}{supp}
\newcommand{\lint}[2]{\underset{#1}{\overset{#2}{{\color{black}\underline{{\color{white}\overline{{\color{black}\int}}\color{black}}}}}}}
\newcommand{\uint}[2]{\underset{#1}{\overset{#2}{{\color{white}\underline{{\color{black}\overline{{\color{black}\int}}\color{black}}}}}}}
\newcommand{\alignint}[2]{\underset{#1}{\overset{#2}{{\color{white}\underline{{\color{white}\overline{{\color{black}\int}}\color{black}}}}}}}
\newcommand{\extint}{\ptxt{ext}\int}
\newcommand{\extalignint}[2]{\ptxt{ext}\alignint{#1}{#2}}
\newcommand{\conv}{\ast}
\newcommand{\pd}[2]{\frac{\partial{}#1}{\partial{}#2}}
\newcommand{\del}{\nabla}
\DeclareMathOperator{\grad}{grad}
\DeclareMathOperator{\curl}{curl}
\let\div\relax
\DeclareMathOperator{\div}{div}
\DeclareMathOperator{\vol}{vol}

% Complex Analysis
\let\Re\relax
\DeclareMathOperator{\Re}{Re}
\let\Im\relax
\DeclareMathOperator{\Im}{Im}
\DeclareMathOperator{\Res}{Res}

% Abstract Algebra
\DeclareMathOperator{\ord}{ord}
\newcommand{\generated}[1]{\left<#1\right>}
\newcommand{\cycle}[1]{\smallPMatrix{#1}}
\newcommand{\id}{\ptxt{id}}
\newcommand{\iso}{\cong}
\DeclareMathOperator{\Aut}{Aut}
\DeclareMathOperator{\SL}{SL}
\DeclareMathOperator{\op}{op}
\newcommand{\isom}[4]{\!\!\!\begin{array}[t]{rcl}#1 & \!\!\!\!\overset{\sim}{\to} & \!\!\!\!#2\\ #3 & \!\!\!\!\mapsto & \!\!\!\!#4\end{array}}
\newcommand{\F}{\mathbb{F}}
\newcommand{\acts}{\;\reflectbox{\rotatebox[origin=c]{-90}{$\circlearrowright$}}\;}
\newcommand{\disjunion}{\mathrel{\text{\scalebox{1.07}{$\perp\mkern-10mu\perp$}}}}
\DeclareMathOperator{\SO}{SO}
\DeclareMathOperator{\stab}{stab}
\DeclareMathOperator{\nullity}{nullity}
\DeclareMathOperator{\Perm}{Perm}
\DeclareMathOperator{\nsubgp}{\vartriangleleft}
\DeclareMathOperator{\notnsubgp}{\ntriangleleft}
\newcommand{\presentation}[2]{\left<#1\;\middle|\;#2\right>}
\DeclareMathOperator{\Char}{char}
\DeclareMathOperator{\fchar}{char}
\DeclareMathOperator{\triv}{triv}
\DeclareMathOperator{\reg}{reg}
\DeclareMathOperator{\std}{std}
\DeclareMathOperator{\Func}{Func}
\DeclareMathOperator{\End}{End}

% Convex Optimization
\let\sectionSymbol\S
\let\S\relax
\newcommand{\S}{\mathbb{S}}
\DeclareMathOperator{\dist}{dist}
\DeclareMathOperator{\dom}{dom}
\DeclareMathOperator{\diag}{diag}
\DeclareMathOperator{\ones}{\mathbbm{1}}
\newcommand{\minimizeOver}[3]{\begin{array}{rl}\underset{#1}{\ptxt{minimize}} & #2\\ \ptxt{subject to} & #3\end{array}}
\newcommand{\maximizeOver}[3]{\begin{array}{rl}\underset{#1}{\ptxt{maximize}} & #2\\ \ptxt{subject to} & #3\end{array}}
\newcommand{\minimizationProblem}[2]{\minimizeOver{}{#1}{#2}}
\newcommand{\maximizationProblem}[2]{\maximizeOver{}{#1}{#2}}
\newcommand{\minimizeOverUnconstrained}[2]{\begin{array}{rl}\underset{#1}{\ptxt{minimize}} & #2\end{array}}
\newcommand{\maximizeOverUnconstrained}[2]{\begin{array}{rl}\underset{#1}{\ptxt{maximize}} & #2\end{array}}
\newcommand{\minimizationUnconstrained}[1]{\minimizeOverUnconstrained{}{#1}}
\newcommand{\maximizationUnconstrained}[1]{\maximizeOverUnconstrained{}{#1}}
\DeclareMathOperator{\argmin}{argmin}
\DeclareMathOperator{\argmax}{argmax}

% Proofs
\newcommand{\st}{s.t.}
\newcommand{\unique}{!}
\newcommand{\iffdef}{\overset{\ptxt{def}}{\Leftrightarrow}}
\newcommand{\eqVertical}{\rotatebox[origin=c]{90}{=}}
\newcommand{\mapsfrom}{\mathrel{\reflectbox{\ensuremath{\mapsto}}}}
\newcommand{\mapsdown}{\rotatebox[origin=c]{-90}{$\mapsto$}\mkern2mu}
\newcommand{\mapsup}{\rotatebox[origin=c]{90}{$\mapsto$}\mkern2mu}
\newcommand{\from}{\!\mathrel{\reflectbox{\ensuremath{\to}}}}
\newcommand{\labeledeq}[1]{\overset{\mathclap{\ptxt{#1}}}{=}}
\newcommand{\eqdef}{\labeledeq{def}}

% Brackets
\newcommand{\paren}[1]{\left(#1\right)}
\renewcommand{\brack}[1]{\left[#1\right]}
\renewcommand{\brace}[1]{\left\{#1\right\}}
\newcommand{\ang}[1]{\left<#1\right>}

% Algorithms
\algrenewcommand{\algorithmiccomment}[1]{\hskip 1em \texttt{// #1}}
\algrenewcommand\algorithmicrequire{\textbf{Input:}}
\algrenewcommand\algorithmicensure{\textbf{Output:}}
\newcommand{\algP}{\ptxt{\textbf{P}}}
\newcommand{\algNP}{\ptxt{\textbf{NP}}}
\newcommand{\algNPC}{\ptxt{\textbf{NP-Complete}}}
\newcommand{\algNPH}{\ptxt{\textbf{NP-Hard}}}
\newcommand{\algEXP}{\ptxt{\textbf{EXP}}}
\DeclareMathOperator{\fl}{fl}

%%%%%%%%%%%%%%%%%%%%%%%%%%%%%
% Other commands
%%%%%%%%%%%%%%%%%%%%%%%%%%%%%
\newcommand{\flag}[1]{\textbf{\textcolor{red}{#1}}}
\let\uSym\u
\let\u\relax
\newcommand{\u}[1]{\underline{#1}}
\let\bSym\b
\let\b\relax
\newcommand{\b}[1]{\textbf{#1}}
\let\iSym\i
\let\i\relax
\newcommand{\i}[1]{\textit{#1}}
\let\scSym\sc
\let\sc\relax
\newcommand{\sc}[1]{\textsc{#1}}

%%%%%%%%%%%%%%%%%%%%%%%%%%%%%%%%%%%%%%%
% Make l's curvy in math environments %
%%%%%%%%%%%%%%%%%%%%%%%%%%%%%%%%%%%%%%%
\mathcode`l="8000
\begingroup
\makeatletter
\lccode`\~=`\l
\DeclareMathSymbol{\lsb@l}{\mathalpha}{letters}{`l}
\lowercase{\gdef~{\ifnum\the\mathgroup=\m@ne \ell \else \lsb@l \fi}}%
\endgroup

%%%%%%%%%%%%%%%%%%%%%%%%%
% Fix \vdots and \ddots %
%%%%%%%%%%%%%%%%%%%%%%%%%
\usepackage{letltxmacro}
\LetLtxMacro\orgvdots\vdots
\LetLtxMacro\orgddots\ddots

\makeatletter
\DeclareRobustCommand\vdots{%
	\mathpalette\@vdots{}%
}
\newcommand*{\@vdots}[2]{%
	% #1: math style
	% #2: unused
	\sbox0{$#1\cdotp\cdotp\cdotp\m@th$}%
	\sbox2{$#1.\m@th$}%
	\vbox{%
		\dimen@=\wd0 %
		\advance\dimen@ -3\ht2 %
		\kern.5\dimen@
		% remove side bearings
		\dimen@=\wd2 %
		\advance\dimen@ -\ht2 %
		\dimen2=\wd0 %
		\advance\dimen2 -\dimen@
		\vbox to \dimen2{%
			\offinterlineskip
			\copy2 \vfill\copy2 \vfill\copy2 %
		}%
	}%
}
\DeclareRobustCommand\ddots{%
	\mathinner{%
		\mathpalette\@ddots{}%
		\mkern\thinmuskip
	}%
}
\newcommand*{\@ddots}[2]{%
	% #1: math style
	% #2: unused
	\sbox0{$#1\cdotp\cdotp\cdotp\m@th$}%
	\sbox2{$#1.\m@th$}%
	\vbox{%
		\dimen@=\wd0 %
		\advance\dimen@ -3\ht2 %
		\kern.5\dimen@
		% remove side bearings
		\dimen@=\wd2 %
		\advance\dimen@ -\ht2 %
		\dimen2=\wd0 %
		\advance\dimen2 -\dimen@
		\vbox to \dimen2{%
			\offinterlineskip
			\hbox{$#1\mathpunct{.}\m@th$}%
			\vfill
			\hbox{$#1\mathpunct{\kern\wd2}\mathpunct{.}\m@th$}%
			\vfill
			\hbox{$#1\mathpunct{\kern\wd2}\mathpunct{\kern\wd2}\mathpunct{.}\m@th$}%
		}%
	}%
}
\makeatother

\newcommand{\B}{
	\begin{tikzpicture}
	\filldraw [fill=red, draw=black] (0, 0) rectangle (0.37, 0.45);
	\draw [line width=0.5mm, white ] (0.1,0.08) -- (0.1,0.38);
	\draw[line width=0.5mm, white ] (0.1, 0.35) .. controls (0.2, 0.35) and (0.4, 0.2625) .. (0.1, 0.225);
	\draw[line width=0.5mm, white ] (0.1, 0.225) .. controls (0.2, 0.225) and (0.4, 0.1625) .. (0.1, 0.1);
	\end{tikzpicture}
}

% Allow custom symbols as arrows in tikz-cd
\tikzset{
  symbol/.style={
    draw=none,
    every to/.append style={
      edge node={node [sloped, allow upside down, auto=false]{$#1$}}}
  }
}
% Example: \arrow[r,symbol=\cong]
% https://tex.stackexchange.com/questions/394154/how-to-include-inclusion-subgroup-relationship-in-tikz-cd-diagram

\author{Professor Alejandro Uribe-Ahumada\\ \small\i{Transcribed by Thomas Cohn}}
\title{Math 591 Lecture 39}
\date{12/16/20} % Can also use \today

\begin{document}
\maketitle
\setlength\RaggedRightParindent{\parindent}
\RaggedRight

\section*{De Rham Cohomology}

\defn{
	Given $M$ a manifold, $k\in\N$ (by our convention, $0\in\N$), we define $Z^{k}(M)=\ker(d:\Omega^{k}\to\Omega^{k+1})$, the set of closed $k$-forms (``cocycles'') and $B^{k}(M)=\im(d:\Omega^{k-1}\to\Omega^{k})$ the set of exact $k$-forms (``coboundaries'').\n
}

\defn{
	The \u{$k$th de Rham group} is the quotient $H^{k}\eqdef{}Z^{k}/B^{k}$. (In our case, this is a quotient vector space, but it can also be defined simply as a quotient group.)\n
}

\defn{
	$\beta_{k}=\dim{}H^{k}(M)$ is the \u{$k$th Betti number} of $M$.\n
}

\par\noindent
Observe: If $k>\dim{}M$, then $\beta_{k}=0$.\n

\par\noindent
Observe: There is a ``compact version'' of this theory, for working with $\Omega_{0}^{k}$, the set of compacly-supported $k$-forms.\n

\ex{
	$H^{0}=\set{f\in{}C^{\infty}\mid{}df=0}$ is the space of locally constant functions. Thus, $\beta_{0}$ is the number of connected components of $M$.\n
}

\ex{
	$M=\R$, $\beta_{0}=1$. What is $\beta_{1}$? Well,\n
	$Z^{1}=\Omega^{1}=\set{f\,dx\mid{}f\in{}C^{\infty}}$.\n
	$B^{1}=\set{dg=g'\,dx\mid{}g\in{}C^{\infty}}$.\n
	Every $f\in{}C^{\infty}$ has an antiderivative, so $H^{1}(\R)=\set{0}$, so $\beta_{1}=0$.\n
}

\prop{
	$H^{k}(\R^{n})\cong\left\{\begin{array}{ll}
		\R & k=0\\
		0 & k>0
	\end{array}\right.$\n
}

\ex{
	$M=\R$. Consider only forms with compactly-supported coefficients ($H^{\cdot}_{C}$). Then\n
	$H^{0}=\set{f\in{}C^{\infty}(\R)\mid{}df=0}=\set{0}$.\n
	To compute $H^{1}$, ask: Which functions $f\in{}C^{\infty}_{0}(\R)$ have anti-derivatives that are of compact support? One can show that this is true iff $\int{}f=0$.\n
}

\par\noindent
Note: $\int:H_{C}^{1}\to\R$ is an isomorphism, so $H_{C}^{1}(\R)\cong\R$.\n

\ex{
	$M=S^{1}$. $M$ is connected, so $\beta_{0}=1$. What about $\beta_{1}$? Well,\n
	$Z^{1}=\Omega^{1}=\set{f\,d\theta\mid{}f:\R\to\R\ptxt{ is $2\pi$-periodic}}$.\n
	$B^{1}=\set{dg=g'\,d\theta\mid{}g\in{}C^{\infty}(S^{1})}$.\n
	Question: Which $2\pi$-periodic functions have $2\pi$-periodic antiderivatives? We can figure this out using Fourier series. We want $f=dg$, so
	\[
		f=\sum_{n\in\Z}a_{n}e^{in\theta}\qquad\Rightarrow\qquad{}g=\frac{1}{i}\sum_{n\in\Z}\frac{a_{n}}{n}e^{in\theta}
	\]
	So we need $a_{0}=0$. In fact, $[f\,d\theta]=[a_{0}\,d\theta]$, where $[\;\cdot\;]$ denotes the cohomology class. Thus, $H^{1}(S^{1})=\R[d\theta]\cong\R$.\n
}

\par\noindent
Observe: This generalizes greatly, to any compact manifold without boundary. It's known as ``Hodge theory''. Fourier series are replaced with the spectral of the Laplacian.\n

\newpage
\section*{General Features}

\begin{itemize}
	\item Covariance: If $F:M\to{}N$ is $C^{\infty}$, we can define, $\forall{}k$, $F^{*}:H^{k}(N)\to{}H^{k}(M)$ by $F^{*}[\omega]=[F^{*}\omega]$ (for $\omega\n{}Z^{k}(N)$). This is well-defined because $F^{*}$ and $d$ commute.
	\item Ring Structure: The wedge product induces a ``cup map'' in cohomology:
	\[
		\map{H^{k}(M)\times{}H^{l}(M)}{H^{k+l}(M)}{([\alpha],[\beta])}{[\alpha\wedge\beta]}
	\]
	Check that this is well-defined:
	\begin{enumerate}[label=\alph*)]
		\item $d(\alpha\wedge\beta)=0$ if $d\alpha=0$ an $d\beta=0$ by the product rule.
		\item $(\alpha+da)\wedge(\beta+db)=\alpha\wedge\beta+\cdots$. The remaining terms are each exact, so it's still the same cohomology class.
	\end{enumerate}
	\item $F^{*}$ is a ring morphism.
	\item $(F\of{}G)^{*}=G^{*}\of{}F^{*}$.
\end{itemize}

\cor{
	Diffeomorphic manifolds have isomorphic cohomology.\n
}

\section*{Homotopy Equivalence}

\defn{
	Two $C^{\infty}$ functions $F,G:M\to{}N$ are \u{homotopic} (to each other) iff there is a smooth map $\Phi:M\times[0,1]\to{}G$ \st{} $\forall{}p\in{}M$, $\Phi(p,0)=F(p)$ and $\Phi(p,1)=G(p)$. $\Phi$ is called a \u{homotopy}.\n
}

\par\noindent
Some notation: $\forall{}t\in\brack{0,1}$, let $\iota_{t}:M\to M\times\brack{0,1}$ where $p\mapsto(p,t)$. Then $F=\Phi\of\iota_{0}$, and $G=\Phi\of\iota_{1}$.\n

\par\noindent
Observe: In Tu's textbook, we use $\R$ instead of $[0,1]$. The two definitions are equivalent, as we can easily extend $\Phi$ from $M\times[0,1]$ to $M\times\R$ with a bump function, and we can simply restrict from $M\times\R$ to $M\times[0,1]$.\n

\thm{
	Being homotopic is an equivalence relation in $C^{\infty}(M,N)$. We write $F\sim{}G$.\n
}

\ex{
	If $X\in\mathfrak{X}(M)$ is complete, then $\forall{}t$, $\phi_{t}$, the flow's time-$t$ map, is homotopic to the identity, $\phi_{0}$.\nn
	Proof: For $t=1$, use $\Phi$. For $t=T$ (nonzero), use $\Phi(p,t)=\phi_{tT}(p)$.\proven
}

\ex{
	Say $\Id,G:\R^{n}\to\R^{n}$, with $G(x)=0$ a constant map. Then $G$ and $I$ are homotopic, with $\Phi(x,t)=(1-t)x$ a homotopy.\n
}

\defn{
	A submanifold $S\subset{}M$ is a \u{deformation retract} of $M$ iff $\exists\Phi:M\times[0,1]\to{}M$ \st{}
	\begin{enumerate}[leftmargin=4\parindent,label=\alph*), itemsep=0pt, topsep=0pt]
		\item $\restr{\Phi}{\set{t=0}}=\Id_{M}$
		\item $\restr{\Phi}{\set{t=1}}$ maps $M$ into $S$.
		\item $\forall t\in[0,1]$, $\forall{}p\in{}S$, $\Phi(p,t)=p$.
	\end{enumerate}\up\n
	Such a $\Phi$ is called a \u{deformation retraction}.\n
}

\ex{
	$M=S^{1}\times(-1,1)$, $S=S^{1}\times\set{0}$. $S$ is a deformation retract of $M$.\n
}

\defn{
	$F:M\to{}N$ is a \u{homotopy equivalence} iff $\exists{}G:N\to{}M$ \st{} $G\of{}F\sim\Id_{M}$ and $F\of{}G\sim\Id_{N}$. If there's a homotopy equivalence $M\to{}N$, we say that $M$ and $N$ are \u{homotopy equivalent}, or that they ``have the same homotopy type''. We say $G$ and $F$ are \u{homotopy inverses}.\n
}

\prop{
	If $S\subset{}M$ is a deformation retract of $M$, then $M$ and $S$ are homotopy equivalent.\nn
	Proof: Let $\Phi:M\times[0,1]\to{}M$ be a deformation retraction of $M$ onto $S$. Define $F:M\to{}S$ to be $F=\restr{\Phi}{M\times\set{1}}^{S}:M\to{}S$, and $G:S\hookrightarrow{}M$ to be the inclusion.\proven
}

\par\noindent
From the perspective of homotopy theory, $\R^{n}$ is the same as a point, the cylinder is the same as the circle. Dimension is not homotopy-invariant!\n

\newpage
\section*{
	Back to de Rham Theory
}

\thm{
	(Homotopy Axiom) Let $F,G:M\to{}N$ be smooth maps, and homotopic to each other. Then $F^{*}=G^{*}$ in cohomology, i.e., $\forall{}k\in\N$, $F^{*},G^{*}:H^{k}(N)\to{}H^{k}(M)$ are equal.\nn
	Proof: First, we'll prove this for $F=I$ , and $G$, the time-$1$ map of a flow. Let $X\in\mathfrak{X}(M)$, and assume it's complete. Say $G=\varphi_{1}$. We need to show $\varphi_{1}^{*}:H^{k}(M)\to{}H^{k}(M)$ is the identity.\nn
	Let $[\omega]\in{}H^{k}(M)$, so $\omega\in{}Z^{k}$ ($d\omega=0$). We need to show $\exists\alpha\in\Omega^{k-1}(M)$ such that $\varphi_{1}^{*}\omega-\omega=d\alpha$.\nn
	Observe: $\frac{d}{dt}\varphi_{t}^{*}\omega=\mathcal{L}_{X}[\varphi_{t}^{*}\omega]=d(\iota_{X}\varphi_{t}^{*}\omega)+\iota_{X}d\varphi_{t}^{*}\omega$. We can commute $d$ and $\varphi_{t}^{*}$, and use the fact that $d\omega=0$, to see that $\frac{d}{dt}\varphi_{t}^{*}\omega=d(\iota_{X}\varphi_{t}^{*}\omega)$. Integrate both sides over $[0,1]$ w.r.t $t$ on each $\bigwedge^{k}T_{p}M$, $\forall{}p\in{}M$, and we obtain
	\[
		\varphi_{1}^{*}\omega-\omega=\int_{0}^{1}\frac{d}{dt}\varphi_{t}^{*}\omega\,dt=\int_{0}^{1}(\iota_{X}\varphi_{t}^{*}\omega)\,dt=d\underbrace{\int_{0}^{1}(\iota_{X}\varphi_{t}^{*}\omega)\,dt}_{\eqdef\alpha}
	\]
	Now, for the general case, say $F,G:M\to{}N$ are homotopic. Then let $\Phi:M\times\R\to{}N$ be a homotopy between them. Let $X=\pd{}{t}$ on $M\times[0,1]$. Then $\varphi_{t}(p,s)=(p,s+t)$. Thus, the following diagram commutes:
	\[
		\begin{tikzcd}[ampersand replacement=\&]
			\&[-5pt] \&[-20pt] M\times\R \arrow[r,"\Phi"] \arrow[dd,"\varphi_{1}"] \& N\\[-15pt]
			M \arrow[urr,hookrightarrow,"\iota_{0}"] \arrow[drr,hookrightarrow,"\iota_{1}",swap] \& \rotatebox[origin=c]{180}{$\circlearrowright$} \\[-15pt]
			\& \& M\times\R \arrow[r,"\Phi"] \& N
		\end{tikzcd}
	\]
	where $\iota_{t}(p)=(p,t)$. So we have $G=\Phi\of\iota_{1}$ and $F=\Phi\of\iota_{0}$.\nn
	Let $[\omega]\in{}H^{k}(N)$.
	\[
		G^{*}[\omega]=(\iota_{1}^{*}\of\Phi^{*})[\omega]=((\varphi_{1}\of\iota_{0})^{*}\of\Phi^{*})[\omega]=(\iota_{0}^{*}\of\varphi_{1}^{*}\of\Phi^{*})[\omega]=\iota_{0}^{*}(\underbrace{\varphi_{1}^{*}(\Phi^{*}[\omega])}_{\substack{=\Phi^{*}[\omega]\\ \ptxt{because }\varphi_{1}^{*}=\Id}})=(\iota_{0}^{*}\of\Phi^{*})[\omega]=F^{*}[\omega]
	\]
	\proven
}

\cor{
	If $M$ and $N$ are homotopy equivalent, then their cohomology is isomorphic.\n
}

\end{document}